% STATUT : À RÉDIGER en dernier.
\chapter{Conclusion}
\label{ch:conclusion}

\textit{[Chapitre à rédiger.]}

% PLAN :
%  1. Ce qui est établi : un mécanisme de dépendance dirigée et à l'ordre
%     entre domaines de contrôle, sa non-transitivité, sa traduction en
%     capital, et la mesure de ce que les cadres existants en expriment.
%  2. Ce qui est démontré comme NON mesurable : la direction, sur données
%     publiques. Avec le test associé (H0 : u_ij = 0 non rejetée).
%  3. Ce qui en découle pour le régulateur : la spécification de registre
%     (survenance horodatée au mois, taxonomie par domaine de contrôle,
%     consolidation des causes communes). C'est la portée du travail
%     au-delà du modèle.
%  4. Limites, sans enrobage : panel d'élicitation restreint, niveau
%     absolu du SCR illustratif, incertitude de queue dominante,
%     paramètres posés listés au chapitre \ref{ch:robustesse}.
%  5. Ouvertures : sous-process P4 sur données client, criticité au sens
%     audit (risque résiduel post-mitigation, méthode ISIS), extension à
%     un portefeuille d'entités.
